
\section{Reálná čísla}


\begin{frame}
\frametitle{Reálná čísla}

\begin{itemize}
\item Celé číslo $X$ v binární soustavě je součet $k$ bitů $b_i$ vynásobených mocninami 2, tedy $ X = \sum_{i=0}^{k-1} b_i 2^{i} $
\item Reálné číslo $X$ v binární soustavě je obdobný součet $k+j$ bitů $b_i$ vynásobených mocninami 2, ale začneme již v záporných mocninách: $X = \sum_{i=-j}^{k} b_i 2^{i}$
\item jistě si všichni pamatujete ze střední školy, že $2^{-j} = \frac{1}{2^j}$
\end{itemize}


\begin{center}
\includegraphics[width=0.65\textwidth]{generated/numbers/float/float_bin.pdf}
\end{center}

\end{frame}


\begin{frame}
\frametitle{Reálná čísla s pevnou desetinnou čárkou}

Čísla s pevnou desetinnou čárkou (fixed point numbers):

\begin{itemize}
\item obdoba čísel s posunutou 0
\item reálné číslo je reprezentované k-bitovým celým číslem se znaménkem, dále zvolíme pevný počet desetinných míst s bitů, kdy $0 \le s \le k$
\item reprezentace -- kód čísla X je $A(X) = \lceil X \cdot 2^s \rceil$
\item rozkódování je funkcí $D(A) = \frac{A}{2^s}$
\item rozsah čísel reprezentovaných je $<-\frac{2^{k-1}}{2^s}, \frac{2^{k-1}-1}{2^s}>$
\item přesnost reprezentace čísel je $\pm \frac{1}{2^s}$.
\end{itemize}

Speciální případ:
\begin{itemize}
\item pokud chcete reprezentovat reálná čísla z intervalu $<0,1)$
\item pak můžete nastavit $s=k$ a použít bezznaménkovou reprezentaci celých čísel
\item dosáhnete lepší přesnosti než odpovídající \texttt{float}, nebo \texttt{double}
\end{itemize}

Některé SIMD instrukce používají pevnou desetinnou čárku.

\end{frame}

\begin{frame}
\frametitle{Reálná čísla s pevnou desetinnou čárkou}

Počítání s čísly s pevnou desetinnou čárkou:

\begin{itemize}
\item součet a rozdíl je součtem a rozdílem celočíselné reprezentace čísel
\item násobení je složitější:
\begin{itemize}
\item $A(X \cdot Y) = (X \cdot Y ) \cdot 2^s = \frac{(X \cdot 2^s)\cdot( Y\cdot 2^s) }{2^s} = \frac{A(X) \cdot A(Y)}{2^s}$
\end{itemize}
\item dělení je obdobné:
\begin{itemize}
\item $A(\frac{X}{Y}) = (\frac{X}{Y}) \cdot 2^s = \frac{(X \cdot 2^s)\cdot2^s }{( Y\cdot 2^s)} = \frac{A(X) \cdot 2^s}{A(Y)}$
\end{itemize}
\item Nejedná se o výrazné zesložitění, protože násobení číslem $2^s$ je posun o s bitů doleva a dělení číslem $2^s$ je posun o s bitů doprava.
\end{itemize}



\end{frame}


\begin{frame}
\frametitle{Plovoucí desetinná čárka}

Plovoucí desetinná čárka (floating point numbers)

\bigskip

Obdoba vědeckého zápisu reálných čísel:
$-123 000 000 000 000.0 = -1.23 \cdot 10^{14} = -1.23\text{E} 14$\\
$-0.000 000 000 000 123 = -1.23 \cdot 10^{-13} = -1.23\text{E}-13$\\

\bigskip

Binární reprezentace používá obdobně mocninu 2:\\
$110 1100 0000 0000.0 = 1.1011 \cdot 2^{14} = 1.1011\text{E} 14 = 29696_{10}$\\
$-0.0000 0000 0000 0001 1101 = -1.1101 \cdot 2^{-16} = -1.1101\text{E}-16 \approx$\\
\phantom{xxxxxxxxxxxxxxxxxxxxxxxxxxxxxxxxxxxxxxxxxx}$\approx0.00002765_{10}$\\

\bigskip

Všimněte si, že každé reálné číslo (kromě 0) začíná v binárním vědeckém zápisu jedničkou.


\end{frame}

\begin{frame}
\frametitle{IEEE-754}

Standard IEEE-754 definuje, jakým způsobem zakódovat reálné číslo do 32, 64 bitů.

Reálné 32-bitové číslo obsahuje:
\begin{itemize}
\item 1 bit znaménko (pozn. existuje 0 i -0)
\item 8 bitů hodnota exponentu s posunutou nulou o 127
\item 23 bitů hodnota mantisy, tedy cifer reprezentujících hodnotu
\end{itemize}
\bigskip
Reálné 64-bitové číslo obsahuje:
\begin{itemize}
\item 1 bit znaménko (pozn. existuje 0 i -0)
\item 11 bitů hodnota exponentu s posunutou nulou o 1023
\item 52 bitů hodnota mantisy, tedy cifer reprezentujících hodnotu
\end{itemize}

\end{frame}


\begin{frame}
\frametitle{IEEE-754}

Příklad:
Reálné číslo $0.828125_{(10)} = 0.5+0.25+0.0625+0.015625=2^{-1}+2^{-2}+2^{-4}+2^{-6} = 0.110101_{(2)}$.\\
Číslo převedeme do vědecké notace: $0.110101 = 1.10101\text{E}-1$.\\
Exponent $e=-1$, reprezentace s posunutou nulou o 127 je $A(-1)=-1+127 = 126$.\\
Abychom neplýtvali zbytečně bity, pro normalizovaná čísla se první 1 mantisy neukládá do binární reprezentace čísla:

\begin{center}
\includegraphics[width=0.4\textwidth]{generated/numbers/float/float_example-cz.pdf}
\end{center}

Zkuste si \url{https://www.h-schmidt.net/FloatConverter/IEEE754.html}
\end{frame}

\begin{frame}
\frametitle{IEEE-754}

Normalizované číslo, je každé číslo, které lze zapsat jako \texttt{1.XXXXX E exp}, kde exp je číslo od -126 do 127, tedy jehož reprezentace je od 1 do 254.

\bigskip
Denormalizovaná čísla jsou čísla, která lze zapsat jako \texttt{0.XXXXX E -126}, tedy například 0.0, nebo všechna čísla v intervalu $(-1.17549\text{E}-38,1.17549\text{E}-38)$, tedy od $(-2^{-126},2^{-126})$.

\bigskip
Pokud je reprezentace exponentu nastavena na samé 1, tedy 255, tzn. hodnota exponentu je 128, pak se jedná o speciální čísla. 
\begin{itemize}
\item pokud je mantisa 0, pak se jedná o nekonečno, podle znaménka \textit{-inf}, nebo \textit{inf}.
\item pokud je mantisa nenulová, pak se jedná o speciální výraz \textit{NaN} -- Not a Number, tedy chybná hodnota čísla, například po výpočtu odmocniny záporného čísla 
\end{itemize}
\end{frame}

\begin{frame}[shrink=5]
\frametitle{IEEE-754}

Přehled reálných čísel:
\begin{tabular}{|c|c|l|}\hline
Exponent & Mantisa &  Hodnota \\ \hline
00000000 & 0 &  0.0 -- čistá nula \\ \hline
00000000 & nenulová &  Denormalizovaná čísla blízká 0 \\ \hline
00000001 & 0 &  nejmenší normalizované číslo se skrytou 1 v mantise \\ \hline
\small 1 až 254 & cokoliv &  normalizovaná čísla, skrytá 1 v mantise \\ \hline
11111111 & 0 &  nekonečno \\ \hline
11111111 & nenulová &  NaN chybná hodnota \\ \hline
\end{tabular}
\bigskip

Denormalizované číslo s nejmenší absolutní hodnotou různé od 0 je:
\begin{itemize}
\item \small exponent = 0 (-126), mantisa=000...0001, hodnota = $2^{-23+(-126)} \approx 1.4\text{E}-45$
\end{itemize}
Normalizované číslo s nejmenší absolutní hodnotou:
\begin{itemize}
\item \small exponent = 1 (-126), mantisa=000...0000, hodnota = $2^{-126} \approx 1.17\text{E}-38$
\end{itemize}
Normalizované číslo s největší absolutní hodnotou:
\begin{itemize}
\item \small exponent = 255 (127), mantisa=111...1111, hodnota = $(2-2^{-23})2^{127} \approx 3.4\text{E}38$
\end{itemize}
\end{frame}


\begin{frame}
\frametitle{IEEE-754 revize 2008 }

Standard IEEE-754 navíc definuje reálné číslo s 16 bity (half precision) a s 128 bity (quad precision).

Reálné 16-bitové číslo obsahuje:
\begin{itemize}
\item 1 bit znaménko 
\item 5 bitů hodnota exponentu s posunutou nulou o 15
\item 10 bitů hodnota mantisy, (plus 11tý bit skrytý neukládaný)
\end{itemize}
\bigskip
Reálné 128-bitové číslo obsahuje:
\begin{itemize}
\item 1 bit znaménko 
\item 15 bitů hodnota exponentu s posunutou nulou o 16383
\item 112 bitů hodnota mantisy, (plus 113 bit skrytý)
\end{itemize}

\end{frame}


\begin{frame}
\frametitle{IEEE-754 }

Porovnání reálných čísel na velikost:
\begin{itemize}
\item Kladná čísla jsou vždy větší než záporná
\item Po odstranění znamének, lze absolutní hodnoty čísel porovnat tak, jak jsou uložené v paměti jako celé číslo bez znaménka
\begin{itemize}
\item To je možné díky zvolené reprezentaci exponentu jako čísla s posunutou nulou
\item Větší exponent větší číslo, při rovnosti exponentů větší mantisa znamená větší číslo
\end{itemize}
\end{itemize}
\end{frame}


\begin{frame}
\frametitle{IEEE-754}

Sčítání nebo odčítání dvou reálných čísel ve vědecké notaci
\begin{enumerate}
\item Převést mantisy čísel na společný exponent, který má hodnotu největšího exponentu (porušíme notaci)
\item Provést součet, nebo rozdíl mantis i se skrytými jedničkami
\item Výsledné číslo normalizovat
\begin{itemize}
\item při sčítání se může exponent zvětšit
\item při odčítání se exponent může zmenšit
\end{itemize}
\end{enumerate}

\end{frame}


\begin{frame}
\frametitle{IEEE-754}

Příklad: Sčítání dvou reálných čísel 31.5+0.75

$31.5_{(10)} = 11111.1_{(2)} = 1.11111\text{E}4$ \phantom{xxx} $0.75_{(10)} = 0.11_{(2)} = 1.1\text{E}-1$

\bigskip
Obě čísla převedeme na stejný exponent 4 a sečteme:\\
\texttt{\phantom{xx}1.11111}\\
\texttt{\phantom{xx}0.000011}\vspace{-6pt}\\
\rule[0pt]{2cm}{0.4pt}\\
\texttt{\phantom{x}10.000001}\\

Výsledné číslo musíme znormalizovat zvýšením exponentu na 5 a tím posunutím desetinné čárky vlevo:\\
$10.000001\text{E}4 = 1.0000001\text{E}5$\\

Toto číslo je reprezentací čísla 32.25.

\end{frame}

\begin{frame}
\frametitle{IEEE-754 -- Násobení}

Násobení dvou reálných čísel:
\begin{enumerate}
\item Exponent výsledku je součet exponentů čísel
\item Mantisa výsledku je součin mantis čísel (i se skrytými jedničkami)
\item Výsledné číslo normalizovat
\begin{itemize}
\item podle výsledku násobení mantis je někdy nutné zvýšit exponent o 1 a vyrotovat mantisu doprava
\end{itemize}
\end{enumerate}
\end{frame}

\begin{frame}
\frametitle{IEEE-754 -- Násobení}

Příklad: Vynásobíme čísla $0.375 \cdot 1.5$

$0.375_{(10)} = 0.011_{(2)} = 1.1\text{E}-2$ \phantom{xxx} $1.5_{(10)} = 1.1_{(2)} = 1.1\text{E}0$

Součin mantis je:\\
\begin{columns}
\begin{column}{0.45\textwidth}
\texttt{\phantom{xxx}11} odpovídá \texttt{1.1}\\
\texttt{\phantom{xx}*11} odpovídá \texttt{1.1}\vspace{-6pt}\\
\rule[0pt]{2cm}{0.4pt}\\
\texttt{\phantom{xxx}11}\\
\texttt{\phantom{xx}11}\vspace{-6pt}\\
\rule[0pt]{2cm}{0.4pt}\\
\texttt{\phantom{x}1001} 

odpovídá \texttt{10.01}, výsledek má 2 desetinná čísla
\end{column}
\hfill
\begin{column}{0.45\textwidth}
\texttt{\phantom{xx}375} odpovídá \texttt{0.375}\\
\texttt{\phantom{xx}*15} odpovídá \texttt{1.5}\vspace{-6pt}\\
\rule[0pt]{2cm}{0.4pt}\\
\texttt{\phantom{x}1875}\\
\texttt{\phantom{x}375}\vspace{-6pt}\\
\rule[0pt]{2cm}{0.4pt}\\
\texttt{\phantom{x}5625} 

odpovídá \texttt{0.5625}, výsledek má 4 desetinná čísla
\end{column}
\end{columns}
\bigskip

Exponent vychází $-2+0=-2$, ale výsledek součinu mantis potřebujeme normalizovat, neboli zvětšit exponent o 1 na $-1$.

Tím dostáváme výsledek $10.01\text{E}-2 = 1.001\text{E}-1$ což je správný výsledek $0.5625_{(10)}$

\end{frame}

\begin{frame}[fragile]
\frametitle{IEEE-754}

Reálná čísla shrnutí:
\begin{itemize}
\item reálná čísla s plovoucí čárkou umí reprezentovat čísla ve velmi velkém rozsahu:
\begin{itemize}
\item float -- absolutní hodnota čísel od $1.175494351\text{E}-38$ do $3.402823466\text{E} + 38$
\item double -- absolutní hodnota čísel od $2.2250738585072014\text{E} - 308$ do $1.7976931348623158\text{E} + 308$
\end{itemize}
\item přesnost čísla se udává na počet validních cifer:
\begin{itemize}
\item float -- v desítkové soustavě 6-7 platných cifer
\item double -- v desítkové soustavě 15-16 platných cifer
\end{itemize}
\end{itemize}

POZOR: Následující \texttt{while} cyklus neskončí:
\begin{minted}{c}
  float a=1.0, step=5e-8;
  while (a*a<1.01) {
    a+=step;
  }
\end{minted}
\end{frame}

